\documentclass{situationist-preprint}

\title{Death Is Bad for the One Who Dies}
\msid{EXI-02-a}
\msfamily{EXI-02}
\msdiscipline{Meaning, Death and the Absurd}
\mstype{Principal manuscript}
\mskeywords{death; deprivation account; Epicurus; Lucretius; symmetry argument; existence requirement; timing problem; Nagel; Feldman}

\begin{document}
\maketitle

\begin{abstract}
Is death bad for the one who dies? Epicurus argued that it is not: when death is present we are not, so there is no one for whom it could be bad. We defend the deprivation account against this argument. On the deprivation account, an event is bad for a subject to the extent that her life as a whole goes worse, by the measure of the goods it contains, than it would have gone had the event not occurred. Death is bad for the one who dies because it deprives her of the goods she would otherwise have had. The key move is to separate two requirements the Epicurean runs together: that there be a subject for whom the event is bad, and that the subject exist at the time of the loss. The first is satisfied, since the subject exists, at the times at which she lives; the second is not a requirement on comparative badness at all, as cases of harm without experience and of loss after the fact show. We then answer the symmetry argument of Lucretius, the problem of locating the time of death's badness, and the objection that the deprivation account proves too much.
\end{abstract}

\section{The problem}

Epicurus wrote to Menoeceus that death is nothing to us, since good and bad lie in sensation and death is the privation of sensation; while we exist death is not present, and when death is present we no longer exist \citep[124--127]{EpicurusMen}. The argument has two premises that later discussion has separated. The \emph{experience requirement} holds that something is bad for a person only if she experiences it as bad. The \emph{existence requirement} holds that something is bad for a person only if she exists when it occurs, or when it is bad for her. Lucretius added a second argument: the time before our birth, an eternity in which we did not exist, is nothing to us, and the time after our death is its mirror image, so it should be nothing to us either \citep[III.832--842, 972--977]{Lucretius}.

The modern response begins with \citet{Nagel1970}, who argued that death is bad, when it is, because of what it deprives us of, and that many goods and evils do not consist in experiences: a man betrayed behind his back is harmed though he never learns of it, and a man reduced by injury to the condition of a contented infant suffers a misfortune though he does not mind it. \citet{Feldman1991,Feldman1992} gave the deprivation account a precise form in terms of the comparison between the values of possible lives, and argued that the badness of death is an eternal fact. \citet{Bradley2009} developed the account further and defended the view that death is bad for its victim at the times after death when she would have been faring well. On the other side, \citet{Rosenbaum1986} defended the Epicurean argument by reformulating it around the experience requirement, and \citet{Silverstein1980} argued that the symmetry problem and the problem of the subject could be met only by a four-dimensional conception of the person's life. The collection of \citet{Fischer1993} gathers the main papers of this debate.

We argue for the deprivation account and against Epicurus. Section~\ref{s:frame} states the account. Section~\ref{s:arg} proves that death is bad for the one who dies whenever it deprives her of a net balance of goods, and shows that the argument requires a subject but not a subject at the time of the loss. Sections~\ref{s:sym} and~\ref{s:obj} answer the symmetry argument and the main objections.

\section{Framework}\label{s:frame}

We follow PID-01 in treating \enquote{person} as a phase sortal of the human organism: a person is an organism in the phase in which it sustains a psychological life, and persists through psychological continuity. Death, in the sense that concerns us, is the end of that phase.

\begin{definition}[Death]\label{d:death}
A person $S$ \emph{dies} at $t$ if $t$ is the last time at which the organism that $S$ is sustains a psychological life, and no psychological life continuous with $S$'s is sustained at any later time.
\end{definition}

Definition~\ref{d:death} leaves open whether the organism persists for a while as a corpse; nothing below depends on it. Its second clause excludes survival. REL-08 argues that survival of death is coherent under a pattern theory of persons but that no evidence available before death could confirm it. Our question concerns death as the end of the person; where a person survives, Definition~\ref{d:death} is not satisfied, and the deprivation account assigns the event its value accordingly.

\begin{definition}[Life value]\label{d:value}
For a possible life $L$ of $S$, $V(L)$ is the sum of the intrinsic values for $S$ of the goods and evils that $L$ contains. The goods include, but are not limited to, pleasant experiences: on the plural account of ETH-01 they include knowledge, achievement and relationships.
\end{definition}

\begin{principle}[Deprivation]\label{p:dep}
The value for $S$ of an event $e$ is $V(L_{@}) - V(L_{\neg e})$, where $L_{@}$ is $S$'s actual life and $L_{\neg e}$ is the life $S$ has in the closest possible world in which $e$ does not occur. The event is \emph{bad for} $S$ if this value is negative, and \emph{good for} $S$ if it is positive.
\end{principle}

Principle~\ref{p:dep} is Feldman's formulation \citep{Feldman1991}. It measures the badness of an event by the difference it makes to the value of the life as a whole, and the lives compared are the whole lives, with all their goods located at whatever times they occur. Three features are important. The account is \emph{comparative}: whether $e$ is bad depends on what the alternative to $e$ would have been. It is \emph{extrinsic}: an event need contain nothing bad in itself to be bad for $S$. And it is \emph{whole-life}: the value it assigns is a fact about $S$'s life taken as a whole.

\begin{assumption}[Tenseless existence]\label{a:tenseless}
Whether a subject stands in a relation is a tenseless matter: $S$ stands in a relation if $S$ exists at some time, and the relation does not require that $S$ exist at a time at which its other relata are located.
\end{assumption}

Assumption~\ref{a:tenseless} is the form that the B-theory of time defended in TSM-01 gives to a familiar point. We stand in relations to things remote from us in time: we are descended from our ancestors, we are admired by our successors, and the author of a book is cited by those who read it after her death. None of these relations requires that its relata coexist at a time.

\section{The argument}\label{s:arg}

\begin{lemma}[Badness without experience]\label{l:exp}
The experience requirement is false: some events are bad for a person though she never experiences them as bad.
\end{lemma}

\begin{proof}
Take Nagel's case of betrayal \citep{Nagel1970}. A woman's friends mock her behind her back and conspire against her interests, though they are scrupulous never to let her know it; her life goes on, as experienced, as before. Her life contains friendships she believes genuine and that are not; by Definition~\ref{d:value}, relationships are among the goods a life contains, and false friendships are not those goods. So her life has less value than the life she would have had without the betrayal, and by Principle~\ref{p:dep} the betrayal is bad for her. She never experiences it as bad. The case of \citet{Nagel1970}'s adult reduced to the mental condition of an infant is the same: he does not mind his condition, and it is a misfortune for him because of the life it has taken from him.
\end{proof}

\begin{lemma}[A subject without a time of loss]\label{l:subject}
If $S$ dies at $t$, there is a subject of whom the comparative fact of Principle~\ref{p:dep} holds, although $S$ does not exist at any time after $t$ at which the goods of $L_{\neg\text{death}}$ would have been located.
\end{lemma}

\begin{proof}
The subject of the comparative fact is $S$, who exists at all the times at which she lives. By Assumption~\ref{a:tenseless}, $S$ can stand in relations whose other relata are located at times at which she does not exist. The comparative fact is such a relation: it relates $S$ to two possible lives, one of which contains goods located after $t$. It is a fact about $S$ and the value of her life as a whole, and it obtains, tenselessly, because $S$ exists and has the life she has. Nothing in Principle~\ref{p:dep} requires that $S$ exist at the times at which the goods she is deprived of would have been located; it requires that those goods belong to the life she would have had.
\end{proof}

\begin{theorem}[Death is bad for the one who dies]\label{t:main}
If $S$ dies at $t$, and the life $S$ would have had had she not died at $t$ contains a greater balance of goods than her actual life, then her death is bad for her. In the ordinary case, in which a person dies while she would otherwise have gone on to a life worth living, this condition is met.
\end{theorem}

\begin{proof}
By Principle~\ref{p:dep}, the value of her death for $S$ is $V(L_{@}) - V(L_{\neg\text{death}})$. Her actual life ends at $t$; the alternative life agrees with it up to $t$ and continues beyond it. By hypothesis the continuation contains a positive balance of goods, so $V(L_{\neg\text{death}}) > V(L_{@})$ and the value of her death for her is negative. By Lemma~\ref{l:subject}, there is a subject of whom this holds, namely $S$; by Lemma~\ref{l:exp}, it does not matter that she never experiences the loss. In the ordinary case the closest world in which she does not die at $t$ is one in which she lives on for a time in roughly the circumstances she was in, and if those circumstances made her life worth living, the continuation has a positive balance.
\end{proof}

\begin{corollary}[Where Epicurus goes wrong]\label{c:epic}
Epicurus's argument fails because deprivation does not require a subject at the time of loss. Its first premise, the experience requirement, is false by Lemma~\ref{l:exp}; its second, read as the claim that a subject must exist at the time at which an event is bad for her, is satisfied in the only sense needed by Lemma~\ref{l:subject}, and is false in the sense the argument requires.
\end{corollary}

\begin{proof}
The argument needs the strong existence requirement: that $S$ exist at the time of the loss. Lemma~\ref{l:subject} shows that the comparative fact constituting the badness of death requires only that $S$ exist, at some time, and not that she exist at the times when the lost goods would have been. So the argument needs a premise that the deprivation account shows to be unmotivated, and the experience requirement on which Epicurus grounded it is false.
\end{proof}

The degree of badness varies. A death at twenty that cuts off decades of a good life is very bad for the one who dies; a death at ninety-five, after which the person would have lived a few months in pain, may be good for her. The account explains both verdicts with one principle. It also explains why deaths that seem equally sudden differ in badness: what matters is not the manner of dying but the life that would otherwise have followed.

\section{The symmetry argument}\label{s:sym}

Lucretius's argument is the harder one. If my death is bad because it deprives me of the goods I would have had had I died later, then my late birth should be bad because it deprives me of the goods I would have had had I been born earlier. We do not regret our late birth; so, by parity, we should not regret our death.

\begin{proposition}[Asymmetry of the alternatives]\label{p:sym}
For a person $S$ who dies at $t$ and began to exist at $t_0$, the closest world in which $S$ dies later than $t$ is, in the ordinary case, much closer than the closest world in which $S$ begins to exist substantially earlier than $t_0$, and the latter often contains no longer good life for $S$.
\end{proposition}

\begin{proof}
A person is an organism in a certain phase (PID-01), and the organism's origin, in a particular pair of gametes, cannot be moved far in time without its ceasing to be the same organism \citep{Kripke1980}. A world in which $S$ begins to exist decades earlier must contain $S$'s parents conceiving her decades earlier, at a time when, in the ordinary case, they were children or not yet born; it differs from the actual world in much of the history that leads to $S$. A world in which $S$ dies later than $t$ need differ from the actual world only in the event that killed her, a car that swerves or a disease that is caught in time. Where an earlier start is possible at all, it is a different life from the beginning, and nothing guarantees that it contains more of $S$'s goods: by Principle~\ref{p:dep}, what is compared is the value of the whole life, not its length.
\end{proof}

Proposition~\ref{p:sym} is the substance of Nagel's reply to Lucretius \citep{Nagel1970}, put in terms of closeness of worlds. It answers the argument without denying the symmetry of the principle. Prenatal nonexistence is bad for $S$ in the rare cases in which there is a close world where she begins earlier and her life is better, and the deprivation account says so. What the symmetry argument assumes is that the alternatives on the two sides are on a par, and Proposition~\ref{p:sym} shows that they are not. A further asymmetry, of attitude, was noted by \citet{BruecknerFischer1986} following \citet[\S 64]{Parfit1984}: we are biased towards future goods and do not care in the same way about goods we have missed. That bias explains why we \emph{feel} differently about the two cases; EXI-05 takes up what it implies for fear. It is not needed for the claim that death is bad.

\section{Objections and replies}\label{s:obj}

\begin{objection}[Timing]\label{o:time}
If death is bad for $S$, it must be bad for her at some time. It cannot be bad for her before she dies, when it has not happened, or after, when she does not exist. So it is bad for her at no time, and therefore not bad for her at all.
\end{objection}

\begin{reply}
The objection assumes that every fact of the form \enquote{$e$ is bad for $S$} must be located at a time. The comparative fact of Principle~\ref{p:dep} is a fact about the value of $S$'s whole life and is not located at any time within it; it is in this sense that \citet{Feldman1992} called the badness of death eternal. \citet{Bradley2009} prefers to say that death is bad for its victim at the times after death at which she would have been well off. Theorem~\ref{t:main} is neutral between these views. On either, the objection's last step fails: from the premise that a fact is not located at a time when $S$ exists it does not follow that it is not a fact about $S$.
\end{reply}

\begin{objection}[Rosenbaum]\label{o:rosen}
The Epicurean can rest the argument on a weaker premise: a state of affairs is bad for a person only if she can experience it \citep{Rosenbaum1986}. Betrayal can be discovered; being dead cannot. The cases of Lemma~\ref{l:exp} therefore do not touch the argument.
\end{objection}

\begin{reply}
The weaker premise draws the line in the wrong place. Suppose the betrayers in Nagel's case could never have been found out, because the woman's circumstances made discovery impossible. Her friendships are no more genuine, and her life no better. Whether a loss could have been discovered bears on whether it could have caused distress, and the deprivation account does not claim that death is bad by causing distress. A principle that ties badness to possible experience is a version of hedonism about well-being, and the goods of Definition~\ref{d:value} that are not experiences are a reason to reject it.
\end{reply}

\begin{objection}[Too much]\label{o:much}
On the deprivation account, everyone who dies at any age is deprived of something, since there is always a close world in which she lives a little longer. So every death is bad, which is incredible for very old or very ill people.
\end{objection}

\begin{reply}
The account does not say that every death is bad. It says that death is bad when the life it ends would otherwise have continued with a positive balance of goods, and the closest world in which a ninety-five-year-old in terminal pain does not die today is one in which she dies next month after a month of pain. Her death is then good for her. The account also allows the badness of death to diminish over a long life, as the goods that would have followed become fewer; EXI-03 argues that in an endless life they eventually give out. Theorem~\ref{t:main} is restricted to the case it states.
\end{reply}

\begin{objection}[Who is deprived]\label{o:who}
At the time of the loss there is at most a corpse. If the corpse is the person (as some animalists hold), the subject exists at the time of loss and Lemma~\ref{l:subject} is idle; if it is not, the deprived subject is the living person, who was never deprived of anything while she lived.
\end{objection}

\begin{reply}
On the view we adopt from PID-01, the corpse, if it is the organism, is not a person, because the phase sortal no longer applies; the goods of Definition~\ref{d:value} are goods of a psychological life, which the corpse does not have. The second horn assumes that being deprived is a state one must be in while living. That is the strong existence requirement again, and Lemma~\ref{l:subject} rejects it: the living person is the subject of the deprivation, and what she is deprived of is located after her life ends.
\end{reply}

\section{What the result does not establish}

The argument does not show that death is the worst thing that can happen to a person, or that it is bad for everyone who dies. It does not settle how much the badness of a death depends on the psychological connections between a person at the time of death and the person she would have become; the time-relative interest account of \citet{McMahan2002}, which discounts the deprivation by the weakness of those connections, is a refinement of Principle~\ref{p:dep} compatible with the result, and it fits the conclusion of PID-06 that connectedness carries most of what matters in survival. It does not show that it is rational to fear death, which is a further question (EXI-05). And it takes no position on survival: if a person survives death in the way REL-08 finds coherent, her death in the biological sense is not death in the sense of Definition~\ref{d:death}.

\section{Conclusion}

Death is bad for the one who dies when it deprives her of the goods she would otherwise have had. The badness is a comparative fact about her life as a whole, and its subject is the person who lived that life. Epicurus's argument fails because it requires that the subject exist at the time of the loss, and deprivation does not require this. Lucretius's argument fails because the alternatives to late birth and to early death are not on a par.

\begin{thebibliography}{99}\small
\bibitem[Bradley(2009)]{Bradley2009} Bradley, B. (2009). \emph{Well-Being and Death}. Oxford University Press.
\bibitem[Brueckner and Fischer(1986)]{BruecknerFischer1986} Brueckner, A.~L. and Fischer, J.~M. (1986). Why is death bad? \emph{Philosophical Studies}, 50(2), 213--221.
\bibitem[Epicurus(Men.)]{EpicurusMen} Epicurus. \emph{Letter to Menoeceus}. In Diogenes Laertius, \emph{Lives of Eminent Philosophers} X.
\bibitem[Feldman(1991)]{Feldman1991} Feldman, F. (1991). Some puzzles about the evil of death. \emph{The Philosophical Review}, 100(2), 205--227.
\bibitem[Feldman(1992)]{Feldman1992} Feldman, F. (1992). \emph{Confrontations with the Reaper}. Oxford University Press.
\bibitem[Fischer(1993)]{Fischer1993} Fischer, J.~M. (ed.) (1993). \emph{The Metaphysics of Death}. Stanford University Press.
\bibitem[Kripke(1980)]{Kripke1980} Kripke, S.~A. (1980). \emph{Naming and Necessity}. Harvard University Press.
\bibitem[Lucretius(DRN)]{Lucretius} Lucretius. \emph{De Rerum Natura}.
\bibitem[McMahan(2002)]{McMahan2002} McMahan, J. (2002). \emph{The Ethics of Killing: Problems at the Margins of Life}. Oxford University Press.
\bibitem[Nagel(1970)]{Nagel1970} Nagel, T. (1970). Death. \emph{No\^us}, 4(1), 73--80.
\bibitem[Parfit(1984)]{Parfit1984} Parfit, D. (1984). \emph{Reasons and Persons}. Oxford University Press.
\bibitem[Rosenbaum(1986)]{Rosenbaum1986} Rosenbaum, S.~E. (1986). How to be dead and not care: A defense of Epicurus. \emph{American Philosophical Quarterly}, 23(2), 217--225.
\bibitem[Silverstein(1980)]{Silverstein1980} Silverstein, H.~S. (1980). The evil of death. \emph{The Journal of Philosophy}, 77(7), 401--424.
\end{thebibliography}

\end{document}
